\section{Step 34 Repeat complete groups and retain the scalar tail}\label{pcm-s34}

Advance all three full pointers by sixteen, consume eight samples and repeat only when another complete group remains. Otherwise fall through to \texttt{mix\_tail}, which returns immediately at zero or processes the final one through seven samples. The scalar tail adds into the same \texttt{v1}. Neither path resets \texttt{v0}. The completed routine now covers short calls, exact blocks and mixed calls under one interface.

\Needspace{12\baselineskip}
\subsection{Write and inspect}
Replace the temporary branch and nop in \texttt{vector\_next}. Keep the surrounding labels and region markers.

\Needspace{10\baselineskip}
\MARSCodeLabel{MIPS assembly}
\begin{minted}[samepage=true]{mars-idea}
    daddiu $a0, $a0, 16
    daddiu $a1, $a1, 16
    daddiu $a2, $a2, 16
    addiu $a3, $a3, -8
    sltiu $t0, $a3, 8
    beq $t0, $zero, mix_block
    nop
\end{minted}
Use recipe \texttt{s34-1}, \texttt{s34-2}, \texttt{s34-3}, \texttt{s34-4} from the companion construction inputs. Pause at \texttt{mix\_ready} before comparing this partial state.

\subsection{Expected observation}
Eight default samples return three clips; nine return four; sixteen return six. A sixty-four-sample call exercises repeated groups without a scalar tail.

\Needspace{11\baselineskip}
\subsection{Check your understanding}
What happens after one group in a nine-sample call?

\begin{itemize}
  \item \textbf{A} The ninth sample is ignored.
  \item \textbf{B} Another sixteen-byte group is read.
  \item \textbf{C} The scalar tail processes one sample.
\end{itemize}
